% =============================================================================
% RAPPORT DE TESTS — Template UPC Projet Informatique L3
% Compiler avec : pdflatex > biber > pdflatex > pdflatex
% =============================================================================
\documentclass[12pt, a4paper]{article}
\usepackage{upc}

% --- Informations du projet --------------------------------------------------
\renewcommand{\typeDocument}{Rapport de Tests}
\renewcommand{\projetNom}{Nom du projet}
\renewcommand{\projetGroupe}{L3Wx}
\renewcommand{\projetEncadrant}{Prénom Nom}
\renewcommand{\projetAuteurs}{Auteur1, Auteur2, Auteur3, Auteur4}
\renewcommand{\projetVersion}{1.0}
\renewcommand{\projetResume}{%
  Ce document présente la campagne de tests menée dans le cadre du développement
  du projet. Il détaille l'ensemble des tests fonctionnels et techniques effectués
  afin de vérifier le bon fonctionnement de l'application, tant du point de vue
  des utilisateurs que des administrateurs.
}

% Commandes utiles pour les tableaux de tests
\newcommand{\valide}{\textcolor{green!60!black}{\textbf{Validé}}}
\newcommand{\echec}{\textcolor{red}{\textbf{Échoué}}}
\newcommand{\partiel}{\textcolor{orange}{\textbf{Partiel}}}

\begin{document}

\makecoverpage
\makesommaire

% =============================================================================
\section{Introduction}
% =============================================================================

Ce document présente la campagne de tests menée dans le cadre du développement
du projet, réalisé dans le cadre de l'UE Projet Informatique de L3 à l'Université
Paris Cité.

Le rapport détaille l'ensemble des tests fonctionnels et techniques effectués
afin de vérifier le bon fonctionnement de l'application, aussi bien du point de
vue des utilisateurs que des administrateurs. Il inclut également une évaluation
des performances du système.

Ces tests ont permis d'identifier les points forts de l'application et les axes
d'amélioration à envisager.

% =============================================================================
\section{Tests généraux}
% =============================================================================

Les tests constituent une étape essentielle du processus de validation du projet.
L'objectif de cette section est de décrire les différents tests réalisés, leurs
objectifs, ainsi que les critères d'acceptation permettant de valider chaque
fonctionnalité.

\subsection{Tests côté utilisateur}

\subsubsection{Création d'un compte}

\noindent\textbf{Objectif :} Vérifier le bon fonctionnement de la création d'un compte utilisateur.\\
\textbf{Pré-requis :} Disposer d'un ordinateur avec accès à internet.

\begin{center}
\begin{tabular}{|c|L{4cm}|L{4.5cm}|c|}
  \hline
  \rowcolor{tabhead}
  \color{white}\textbf{N°} &
  \color{white}\textbf{Action} &
  \color{white}\textbf{Attendu} &
  \color{white}\textbf{Résultat} \\
  \hline
  UA1 & Cliquer sur ``Créer un compte'' &
        Ouverture de la page d'inscription & \valide \\
  \hline
  UA2 & Saisir une adresse e-mail invalide &
        Message d'erreur de format & \valide \\
  \hline
  UA3 & Saisir un e-mail valide sans mot de passe &
        Message ``Veuillez saisir un mot de passe'' & \valide \\
  \hline
  UA4 & Saisir un e-mail valide et un mot de passe trop court &
        Message d'erreur de mot de passe & \valide \\
  \hline
  UA5 & Saisir un e-mail valide et un mot de passe valide &
        Création du compte et connexion automatique & \valide \\
  \hline
  UA6 & Saisir un e-mail déjà utilisé &
        Message ``Compte déjà existant'' & \valide \\
  \hline
\end{tabular}
\end{center}

\subsubsection{Connexion}

\noindent\textbf{Objectif :} Vérifier le bon fonctionnement de la connexion.\\
\textbf{Pré-requis :} Posséder un compte existant.

\begin{center}
\begin{tabular}{|c|L{4cm}|L{4.5cm}|c|}
  \hline
  \rowcolor{tabhead}
  \color{white}\textbf{N°} &
  \color{white}\textbf{Action} &
  \color{white}\textbf{Attendu} &
  \color{white}\textbf{Résultat} \\
  \hline
  UB1 & Saisir des identifiants corrects &
        Connexion et redirection vers la page principale & \valide \\
  \hline
  UB2 & Saisir un mot de passe incorrect &
        Message d'erreur d'authentification & \valide \\
  \hline
  UB3 & Laisser les champs vides &
        Message ``Champs requis'' & \valide \\
  \hline
\end{tabular}
\end{center}

\subsubsection{Fonctionnalité principale}

\noindent\textbf{Objectif :} Vérifier le bon fonctionnement de [fonctionnalité].\\
\textbf{Pré-requis :} Être connecté en tant qu'utilisateur.

\begin{center}
\begin{tabular}{|c|L{4cm}|L{4.5cm}|c|}
  \hline
  \rowcolor{tabhead}
  \color{white}\textbf{N°} &
  \color{white}\textbf{Action} &
  \color{white}\textbf{Attendu} &
  \color{white}\textbf{Résultat} \\
  \hline
  UC1 & [Action 1] & [Résultat attendu] & \valide \\
  \hline
  UC2 & [Action 2] & [Résultat attendu] & \valide \\
  \hline
  UC3 & [Action 3] & [Résultat attendu] & \partiel \\
  \hline
\end{tabular}
\end{center}

\subsection{Tests côté administrateur}

\subsubsection{Gestion des utilisateurs}

\noindent\textbf{Objectif :} Vérifier les fonctionnalités d'administration des comptes.\\
\textbf{Pré-requis :} Être connecté en tant qu'administrateur.

\begin{center}
\begin{tabular}{|c|L{4cm}|L{4.5cm}|c|}
  \hline
  \rowcolor{tabhead}
  \color{white}\textbf{N°} &
  \color{white}\textbf{Action} &
  \color{white}\textbf{Attendu} &
  \color{white}\textbf{Résultat} \\
  \hline
  AA1 & Accéder à la liste des utilisateurs &
        Affichage de tous les comptes & \valide \\
  \hline
  AA2 & Modifier le rôle d'un utilisateur &
        Modification enregistrée & \valide \\
  \hline
  AA3 & Supprimer un compte utilisateur &
        Compte supprimé, confirmation affichée & \valide \\
  \hline
\end{tabular}
\end{center}

\subsubsection{Gestion du contenu}

\noindent\textbf{Objectif :} Vérifier la gestion du contenu par l'administrateur.

\begin{center}
\begin{tabular}{|c|L{4cm}|L{4.5cm}|c|}
  \hline
  \rowcolor{tabhead}
  \color{white}\textbf{N°} &
  \color{white}\textbf{Action} &
  \color{white}\textbf{Attendu} &
  \color{white}\textbf{Résultat} \\
  \hline
  AB1 & Ajouter un nouvel élément de contenu &
        Élément ajouté et visible & \valide \\
  \hline
  AB2 & Modifier un élément existant &
        Modifications sauvegardées & \valide \\
  \hline
  AB3 & Supprimer un élément &
        Élément supprimé et inaccessible & \valide \\
  \hline
\end{tabular}
\end{center}

\subsection{Tests de performance}

\noindent\textbf{Objectif :} Évaluer les performances de l'application sous charge.

\begin{center}
\begin{tabular}{|L{4cm}|c|c|c|}
  \hline
  \rowcolor{tabhead}
  \color{white}\textbf{Scénario} &
  \color{white}\textbf{Utilisateurs} &
  \color{white}\textbf{Temps moyen} &
  \color{white}\textbf{Résultat} \\
  \hline
  Chargement page accueil & 1   & < 1\,s   & \valide \\
  \hline
  Chargement page accueil & 10  & < 2\,s   & \valide \\
  \hline
  Requête API principale  & 1   & < 3\,s   & \valide \\
  \hline
  Requête API principale  & 10  & < 5\,s   & \partiel \\
  \hline
\end{tabular}
\end{center}

% =============================================================================
\section{Tests qualitatifs}
% =============================================================================

[Décrire ici des tests qualitatifs spécifiques à votre application : qualité
des réponses générées, pertinence du contenu, etc.]

\subsection{Objectif}

[Décrire l'objectif de ces tests qualitatifs.]

\subsection{Structure de test}

Chaque test est structuré comme suit :
\begin{itemize}
  \item \textbf{Entrée :} [description de l'entrée testée]
  \item \textbf{Sortie attendue :} [critère de qualité]
  \item \textbf{Évaluation :} [méthode d'évaluation]
\end{itemize}

\subsection{Résultats}

[Présenter les résultats des tests qualitatifs sous forme de tableau ou
de description.]

\begin{center}
\begin{tabular}{|c|L{5cm}|c|c|}
  \hline
  \rowcolor{tabhead}
  \color{white}\textbf{N°} &
  \color{white}\textbf{Cas de test} &
  \color{white}\textbf{Score} &
  \color{white}\textbf{Résultat} \\
  \hline
  Q1 & [Description du cas de test 1] & [4/5] & \valide \\
  \hline
  Q2 & [Description du cas de test 2] & [3/5] & \partiel \\
  \hline
  Q3 & [Description du cas de test 3] & [5/5] & \valide \\
  \hline
\end{tabular}
\end{center}

% =============================================================================
\section{Conclusion}
% =============================================================================

[Synthétiser les résultats des tests : nombre de tests réussis/échoués,
points forts et points d'amélioration identifiés.]

Les tests réalisés ont permis de valider le bon fonctionnement de l'application
dans la majorité des cas. Les points d'amélioration identifiés sont les suivants :

\begin{itemize}
  \item [Point d'amélioration 1]
  \item [Point d'amélioration 2]
\end{itemize}

\end{document}
